\documentclass[10pt,a4paper]{amsart}
\usepackage[margin=19mm]{geometry}
\usepackage{amsmath,amssymb,mathtools}
\usepackage[scr=rsfs,cal=euler]{mathalfa}
\usepackage{xfrac}
\usepackage[unicode,hidelinks,bookmarks=false]{hyperref}
\newcommand{\ie}{i.e.\ }
\newcommand{\cf}{cf.\ }
\newcommand{\resp}{respectively\ }
\newcommand{\Subsection}{\subsection}
\providecommand{\nobreakdash}{\nobreak\hbox{-}}
\setlength{\emergencystretch}{2em}
\multlinegap0pt
\usepackage{amssymb}
\usepackage{tikz}
\usepackage{mathtools}
\usepackage{xcolor}
\usepackage[normalem]{ulem}
\newtheorem{proposition}{Proposition}
\newcommand{\upbar}[1]{\,\overline{\! #1}}
\title{Freidlin-Wentzell solutions of discrete Hamilton-Jacobi equations}
\author{Michele Aleandri, Davide Gabrielli, Giulia Pallotta}
\date{Source: arXiv:2501.12505v2 (CC BY 4.0)}
\begin{document}
\maketitle

\noindent\emph{Source and licence.} Freidlin-Wentzell solutions of discrete Hamilton-Jacobi equations. Version arXiv:2501.12505v2 (CC BY 4.0), distributed by arXiv under the Creative Commons Attribution 4.0 International licence, http://creativecommons.org/licenses/by/4.0/, as recorded in the arXiv licence metadata for this exact version and shown on the abstract page https://arxiv.org/abs/2501.12505v2. Full licence notice: \texttt{licenses/arxiv-2501.12505-CC-BY-4.0.txt}.\par\smallskip

\noindent\textit{Editorial scope.} Counted unit: the article's proposition propFW (source0.tex lines 620--622). The article's complete original proof of it is delivered with it (source0.tex lines 623--645, one \texttt{proof} environment of 23 lines). No mathematics is written, repaired or re-derived: the statement and the proof are the article's own lines, in the article's own notation, and no retained line is substituted or re-flowed.  Retained uncounted as the source-local context the proof needs: lines 208--210; lines 397--399.  Nothing was omitted from any retained span, so the package carries no omission record.  No retained line contains a citation, so the wrapper carries no bibliography and no bibliography entry.  No retained line contains a figure environment, an image-inclusion command or drawing code, so no image is delivered and the package declares no asset dependency.  Editorial formatting only, in addition: an AMS \texttt{amsart} preamble that restates the article's own declarations and macros used by the retained lines, namely the environments \texttt{proposition} on separate counters, together with the standard packages those lines need.  The retained environments print in this document as Proposition 1 in order of appearance, whereas the article prints the same items with the numbering of its own full document.  The complete native source of the article as distributed in the arXiv e-print for this exact version is bundled unchanged as \texttt{sources/171935/source0.tex}.
	\begin{equation}\label{1HJ}
		W(x)+\min_{y:(x,y)\in E}\left\{\Delta(x,y)\right\}=\min_{y:(y,x)\in E}\left\{W(y)+\Delta(y,x)\right\}\,,\qquad \forall x\in V\,.
	\end{equation}

	\begin{equation} \label{FW solution}
		\mathscr{FW}(x):=\lim_{N\to +\infty}-\frac 1N\log \pi_N(x)=\inf_{\tau\in \mathcal T_x} \Delta(\tau)-\inf_{\tau\in \mathcal T} \Delta(\tau)\,, \qquad x\in V\,.
	\end{equation}

\begin{proposition}\label{propFW}
	The $FW$-quasipotential in \eqref{FW solution} is always an element of $\mathcal W_0$, the set of normalized solutions of \eqref{1HJ}.
\end{proposition}

\begin{proof}
	The validity of \eqref{1HJ} is checked for each node. Fix $x\in V$ and let $y^*\in V$ be such that $(x,y^*)\in E_0$ and $\tau^*\in \mathcal T_x$ be the minimal arborescence in \eqref{FW solution}. Then $\Delta(x,y^*)=0$ and $\Delta(\tau^*)$ is minimal in $\mathcal T_x$. 
	
	From the setup,  $\gamma^*:=(x,y^*)\cup \tau^* \in  \mathcal{U}^{\mathrm{in}}_x$ and moreover 
	\begin{equation}\label{fox}
		\Delta(\gamma^*)=\min_{\gamma\in \mathcal{U}^{\mathrm{in}}_x} \Delta(\gamma)=\Delta(\tau^*)\,.
	\end{equation} 
	In fact,  any $\gamma\in \mathcal{U}^{\mathrm{in}}_x$ can be written as $\gamma=(x,y)\cup \tau$ for a suitable $y\in V$ and $\tau\in \mathcal T_x$ and $\Delta(\gamma)=\Delta(x,y)+\Delta (\tau)$. Since both $\Delta(x,y^*)$ and $\Delta (\tau^*)$ are minimal (respectively on $(x,y)\in E$ and $\mathcal T_x$) we  proved \eqref{fox}.
	
	For any $y\in V$ , up to a constant, $\mathscr{FW}(y)=\Delta(\tau)$ for a suitable $\tau\in \mathcal T_y$ and moreover when $(y,x)\in E$ the graph $(y,x)\cup \tau\in \mathcal{U}^{\mathrm{in}}_x$. By \eqref{fox} 
	\begin{equation}\label{fox2}
		\mathscr{FW}(y)+\Delta(y,x)=\Delta((y,x)\cup \tau)\geq \Delta(\gamma^*)=\mathscr{FW}(x)\,,\qquad \forall (y,x)\in E\,,
	\end{equation}
	that means $\mathscr{FW}\in \underline{\mathcal W}$.
	It remains to prove that there exists an $(y',x)\in E$ such that the equality holds in \eqref{fox2}, i.e. $\mathscr{FW}\in \upbar{\mathcal W}$. Consider $(y',x)$ the unique edge entering in $x$ and belonging to the unique cycle in
	$\gamma^*\in \mathcal{U}^{\mathrm{in}}_x$ (recall that $\gamma^*$ is carachterized by \eqref{fox}). 
	
	Since $\gamma^*\setminus (y',x)\in \mathcal T_{y'}$
	\begin{equation}
		\mathscr{FW}(y')+\Delta(y',x)\leq \Delta(\gamma^*\setminus (y',x))+\Delta(y',x)=\Delta(\gamma^*)=\mathscr{FW}(x)\,.
	\end{equation}
	Last inequality  together with the inequality \eqref{fox2}, implies the equality.
\end{proof}

\end{document}
